\documentclass[11pt]{article}
\usepackage[margin=1in]{geometry}
\usepackage[T1]{fontenc}
\usepackage[utf8]{inputenc}
\usepackage{lmodern}
\usepackage{microtype}
\usepackage[table]{xcolor}
\usepackage{amsmath,amssymb,amsfonts,mathtools,bm}
\IfFileExists{physics.sty}{%
  \usepackage{physics}%
}{%
  % Portable subset used by this project when the optional physics package
  % is not installed in the local TeX distribution.
  \NewDocumentCommand{\pdv}{o m m}{%
    \IfNoValueTF{##1}{\frac{\partial ##2}{\partial ##3}}%
    {\frac{\partial^{##1} ##2}{\partial ##3^{##1}}}%
  }
  \NewDocumentCommand{\dv}{o m m}{%
    \IfNoValueTF{##1}{\frac{\mathrm{d} ##2}{\mathrm{d} ##3}}%
    {\frac{\mathrm{d}^{##1} ##2}{\mathrm{d} ##3^{##1}}}%
  }
  \providecommand{\dd}{\mathop{}\!\mathrm{d}}
  \providecommand{\grad}{\nabla}
  \providecommand{\abs}[1]{\left\lvert ##1\right\rvert}
  \providecommand{\norm}[1]{\left\lVert ##1\right\rVert}
  \providecommand{\ket}[1]{\left\lvert ##1\right\rangle}
  \providecommand{\bra}[1]{\left\langle ##1\right\rvert}
}
\usepackage{booktabs,array,longtable,tabularx}
\usepackage{hyperref}
\usepackage{enumitem}
\usepackage{fancyhdr}
\usepackage{titlesec}
\usepackage{tcolorbox}
\usepackage{ragged2e}
\usepackage{setspace}
\IfFileExists{siunitx.sty}{%
  \usepackage{siunitx}%
}{%
  % Minimal text-safe fallback for the small number of \SI and \si calls in
  % this repository. Full siunitx formatting is used when available.
  \providecommand{\si}[1]{\ensuremath{\mathrm{##1}}}
  \providecommand{\SI}[2]{\ensuremath{##1\,\mathrm{##2}}}
}
\hypersetup{colorlinks=true, linkcolor=blue!50!black, urlcolor=blue!50!black, citecolor=blue!50!black}
\definecolor{TitleBlue}{HTML}{163A5F}
\definecolor{AccentTeal}{HTML}{2D6F73}
\definecolor{SoftBlue}{HTML}{EAF3F8}
\definecolor{SoftTeal}{HTML}{EDF8F6}
\definecolor{SoftGray}{HTML}{F5F7FA}
\definecolor{RuleGray}{HTML}{D7DEE8}
\pagestyle{fancy}
\fancyhf{}
\lhead{RadiationEffects\_proj2}
\rhead{Technical Notes}
\cfoot{\thepage}
\setlength{\headheight}{14pt}
\setlength{\parskip}{0.5em}
\setlength{\parindent}{0pt}
\onehalfspacing
\numberwithin{equation}{section}
\titleformat{\section}{\Large\bfseries\color{TitleBlue}}{\thesection}{0.6em}{}
\titleformat{\subsection}{\large\bfseries\color{AccentTeal}}{\thesubsection}{0.6em}{}
\titleformat{\subsubsection}{\normalsize\bfseries\color{TitleBlue}}{\thesubsubsection}{0.6em}{}
\newtcolorbox{overviewbox}{colback=SoftBlue,colframe=TitleBlue,boxrule=0.8pt,arc=2mm,left=3mm,right=3mm,top=2mm,bottom=2mm}
\newtcolorbox{physicsbox}{colback=SoftTeal,colframe=AccentTeal,boxrule=0.8pt,arc=2mm,left=3mm,right=3mm,top=2mm,bottom=2mm}
\newtcolorbox{notebox}{colback=SoftGray,colframe=AccentTeal,boxrule=0.6pt,arc=2mm,left=3mm,right=3mm,top=2mm,bottom=2mm}
\newcommand{\DocTitle}[3]{%
\begin{center}
{\LARGE \bfseries \color{TitleBlue} #1}\\[0.5em]
{\large \color{AccentTeal} #2}\\[0.8em]
{\normalsize #3}
\end{center}
\vspace{0.5em}}
\newenvironment{symtable}{\rowcolors{2}{SoftGray}{white}\begin{longtable}{>{\RaggedRight\arraybackslash}p{0.18\textwidth} >{\RaggedRight\arraybackslash}p{0.25\textwidth} >{\RaggedRight\arraybackslash}p{0.47\textwidth}}\toprule \textbf{Symbol / acronym} & \textbf{Name} & \textbf{Meaning in this document} \\ \midrule\endfirsthead\toprule \textbf{Symbol / acronym} & \textbf{Name} & \textbf{Meaning in this document} \\ \midrule\endhead}{\bottomrule\end{longtable}}


\newcommand{\ProjectTOC}{%
\vspace{0.6em}
{\hypersetup{linkcolor=TitleBlue,urlcolor=TitleBlue,citecolor=TitleBlue}\tableofcontents}
\vspace{1.0em}}

\newcommand{\SymbolsAcronymsSection}{%
\section*{Symbols and acronyms used in this document}%
\addcontentsline{toc}{section}{Symbols and acronyms used in this document}}

\newcommand{\rvec}{\bm{r}}
\newcommand{\qvec}{\bm{q}}
\newcommand{\nvec}{\bm{n}}
\newcommand{\xvec}{\bm{x}}
\newcommand{\kB}{k_{\mathrm{B}}}
\newcommand{\qp}{\mathrm{qp}}
\newcommand{\JJ}{\mathrm{JJ}}
\newcommand{\mfp}{\Lambda}
\allowdisplaybreaks

\usepackage{parskip}
\rhead{Slides Outline}

\begin{document}

\DocTitle{RadiationEffects\_proj2 Slide Outline}{Project overview roadmap + deep-dive lecture outline for Module 4}{Planning document for building a physics-lecture-style presentation from \texttt{project\_plan.pdf} and \texttt{module4\_2d\_ballistic\_diffusive\_thermal\_transport.pdf}}

\begin{overviewbox}
\textbf{Purpose.} This document is not the final PowerPoint deck. It is a slide-planning outline for \textbf{Radiation Effects Project 2}. The first part now follows a tighter structure: \textbf{one project-overview slide plus one slide per module} (Modules 1--7). The second part remains a \textbf{10-slide deep dive into Module 4}. Each slide entry is written to help the eventual talk feel like a physics lecture rather than a project-management briefing.
\end{overviewbox}

\ProjectTOC

\SymbolsAcronymsSection
\begin{symtable}
BD & Ballistic-diffusive & Reduced thermal-transport model that treats boundary-originating carriers ballistically and internally scattered carriers diffusively.\\
BTE & Boltzmann transport equation & Microscopic transport equation for the carrier distribution function.\\
FEM & Finite element method & Planned numerical discretization framework for later geometric implementations.\\
RTA & Relaxation-time approximation & Collision model in which the distribution relaxes toward local equilibrium over a characteristic time \(\tau\).\\
\(C_j\) & Defect concentration & Concentration of defect species \(j\) used in Module 3 and later coupled modules.\\
\(\rho\) & Space-charge density & Net charge density entering Poisson's equation in Module 2.\\
\(\phi\) & Electric potential & Scalar electrostatic potential solved in Module 2.\\
\(n,p\) & Electron and hole densities & Carrier populations used in Module 5 drift-diffusion transport.\\
\(\qvec_b\) & Ballistic heat flux & Boundary-memory heat-flux contribution retained in Module 4.\\
\(u_m\) & Medium/scattered energy & Internal-energy field associated with the scattered part of the carrier population.\\
\(Q\) & Volumetric heat source & Radiation-related, electrical, or externally imposed source term for Module 4.\\
\(\mfp\) & Mean free path & Average distance traveled by a heat carrier between scattering events.\\
\(\mathrm{Kn}\) & Knudsen number & Ratio \(\mathrm{Kn}=\mfp/L\), used to organize diffusive versus ballistic transport regimes.\\
\end{symtable}

\section{How to use this slide outline}
Each slide entry gives six items:
\begin{enumerate}[leftmargin=1.6em]
    \item the main question the slide should answer,
    \item the core message the audience should remember,
    \item the equation or equations that should be shown,
    \item a figure idea,
    \item a physics drawing idea,
    \item the speaker emphasis that should guide the verbal explanation.
\end{enumerate}
The final slides should stay visually sparse. The derivation, physical explanation, and comparisons should mostly live in the speaker notes and in the spoken delivery.

\begin{notebox}
\textbf{Presentation principle.} The overview part should feel like a physics-guided narrative of why each module exists, not like a software inventory. The Module 4 part should not begin with the final reduced PDE. It should begin with the physical reason Fourier conduction is insufficient, then move from microscopic heat carriers to the ballistic-diffusive split, and only then show the final 2D target model used in \texttt{RadiationEffects\_proj2}.
\end{notebox}

\section{Part I: Project overview slides (8 slides total)}

\section{Slide 1: Project overview --- the full physics chain}
\begin{physicsbox}
\textbf{Main question.} What is the full physical story that this project is trying to model?
\end{physicsbox}

\textbf{Core message.} The project is built to explain how a radiation insult in silicon evolves into material damage, altered thermal and electrical fields, modified carrier transport, and finally degraded device behavior.

\textbf{Equations to show:}
\begin{equation*}
\dot{q}_{\mathrm{rad}}(x,y,t)
\;\longrightarrow\;
C_j(x,y,t)
\;\longrightarrow\;
\rho(x,y,t),\, \phi(x,y,t),\, T(x,y,t)
\;\longrightarrow\;
n,p,J_n,J_p
\;\longrightarrow\;
\text{device degradation}.
\end{equation*}

\textbf{Figure idea.} One end-to-end roadmap graphic with the modules laid along the physical chain from radiation deposition to device metrics.

\textbf{Physics drawing idea.} A silicon device cross section with an incident particle track, a localized damage region, charged defects, a thermal hot spot, bent field lines, and distorted current flow.

\textbf{Speaker emphasis.} Set the tone immediately: this is a mechanism-resolved physics framework. The point is not just to simulate radiation deposition, but to preserve the downstream physics that determines whether the device actually degrades.

\section{Slide 2: Module 1 --- radiation deposition and damage seeding}
\begin{physicsbox}
\textbf{Main question.} What does Module 1 provide to the rest of the framework?
\end{physicsbox}

\textbf{Core message.} Module 1 uses radiation transport to generate the initial spatial forcing: where energy is deposited, what interactions occur, and how that deposition seeds the later defect-evolution problem.

\textbf{Equations to show:}
\begin{align*}
S_{\mathrm{dep}}(x,y,t) &\equiv \text{deposited-energy or source map},\\
\Phi &\longrightarrow S_{\mathrm{dep}},\; \text{interaction tallies},\; \text{damage precursors}.
\end{align*}

\textbf{Figure idea.} A deposition map from a particle-transport run, with track structure or localized hotspots highlighted.

\textbf{Physics drawing idea.} Incident particle entering silicon, showing scattering events, energy transfer, and a map of where the material first receives the radiation insult.

\textbf{Speaker emphasis.} Explain that Module 1 is the interface between external radiation physics and internal material/device physics. It does not yet tell us the evolved damage state; it gives the initial spatial insult.

\section{Slide 3: Module 2 --- electrostatics with defect-dependent space charge}
\begin{physicsbox}
\textbf{Main question.} How does a damaged defect population alter the electrostatic state of the device?
\end{physicsbox}

\textbf{Core message.} Module 2 translates the current charged-defect and dopant populations into electric potential and field maps through Poisson's equation.

\textbf{Equations to show:}
\begin{align*}
\nabla\cdot(\epsilon\nabla\phi) &= -\rho,\\
\rho &= q\left(p-n+N_D^+-N_A^-+\sum_j z_j C_j\right),\\
\bm{E} &= -\nabla\phi.
\end{align*}

\textbf{Figure idea.} A 2D potential map and electric-field-line plot over a device cross section with a localized defect-rich region.

\textbf{Physics drawing idea.} Show charged defects locally bending field lines and perturbing depletion structure or junction fields.

\textbf{Speaker emphasis.} Make clear that Module 2 does not evolve the defects; it translates the current defect state into an electrostatic consequence that later feeds transport.

\section{Slide 4: Module 3 --- defect evolution in space and time}
\begin{physicsbox}
\textbf{Main question.} How does the radiation-created defect population evolve after the initial event?
\end{physicsbox}

\textbf{Core message.} Module 3 advances defect concentrations through diffusion, reaction, clustering, dissociation, and annealing, with coefficients that can depend on temperature, fields, and local material state.

\textbf{Equations to show:}
\begin{equation*}
\frac{\partial C_j}{\partial t}
=\nabla\cdot\bigl(D_j\nabla C_j\bigr)
+R_j(C_1,C_2,\dots,T,E).
\end{equation*}
\begin{align*}
D_j(T) &= D_{0,j}\exp\!\left(-\frac{E_{D,j}}{\kB T}\right),\\
k_j(T) &= k_{0,j}\exp\!\left(-\frac{E_{A,j}}{\kB T}\right).
\end{align*}

\textbf{Figure idea.} A time sequence of defect concentration maps spreading, reacting, and relaxing.

\textbf{Physics drawing idea.} Vacancies and interstitials migrating, recombining, and annealing in a silicon lattice, with arrows indicating competing timescales.

\textbf{Speaker emphasis.} Stress that Module 3 is the engine that converts the initial deposition pattern into a changing damage landscape. Without this step, the project would assume the material damage stays frozen.

\section{Slide 5: Module 4 --- ballistic-diffusive thermal transport}
\begin{physicsbox}
\textbf{Main question.} Why is thermal transport its own module, and what makes it different from a plain heat equation?
\end{physicsbox}

\textbf{Core message.} Module 4 models heat transport in a regime where finite carrier flight lengths and boundary memory matter, so the thermal problem is richer than ordinary Fourier conduction.

\textbf{Equations to show:}
\begin{equation*}
\tau\frac{\partial^2 u_m}{\partial t^2}+\frac{\partial u_m}{\partial t}
=\nabla\cdot\left(\frac{k}{C}\nabla u_m\right)-\nabla\cdot\qvec_b
+\left(Q+\tau\frac{\partial Q}{\partial t}\right).
\end{equation*}

\textbf{Figure idea.} A 2D thermal field with ballistic boundary rays and a smoother scattered-temperature background.

\textbf{Physics drawing idea.} Localized hot spot, outward-moving heat carriers, partial scattering, and boundary-emitted carriers contributing nonlocal thermal memory.

\textbf{Speaker emphasis.} Keep this slide high-level. The goal is to show why Module 4 exists in the chain and to set up the later deep dive.

\section{Slide 6: Module 5 --- carrier transport and recombination under damage}
\begin{physicsbox}
\textbf{Main question.} Once fields, temperature, and defects are known, how is device behavior predicted?
\end{physicsbox}

\textbf{Core message.} Module 5 converts the damaged material state into an electrical-performance problem by solving carrier drift, diffusion, and recombination in the perturbed field and trap landscape.

\textbf{Equations to show:}
\begin{align*}
J_n &= q\mu_n nE + qD_n\nabla n,\\
J_p &= q\mu_p pE - qD_p\nabla p,\\
\frac{\partial n}{\partial t} &= \frac{1}{q}\nabla\cdot J_n + G - R,\\
\frac{\partial p}{\partial t} &= -\frac{1}{q}\nabla\cdot J_p + G - R.
\end{align*}

\textbf{Figure idea.} A current-density map flowing around a damaged region, with local recombination or trapping zones indicated.

\textbf{Physics drawing idea.} Electrons and holes drifting under a bent field while defects act as recombination centers and mobility modifiers.

\textbf{Speaker emphasis.} This is where the materials physics becomes a device-performance prediction. Emphasize that damage matters electrically because it modifies both the driving fields and the recombination/trapping landscape.

\section{Slide 7: Module 6 --- self-consistent multiphysics coupling}
\begin{physicsbox}
\textbf{Main question.} Why can the modules not remain independent forever?
\end{physicsbox}

\textbf{Core message.} In the final framework, defects, temperature, fields, and carrier populations all feed back on one another, so the full problem must be iterated to consistency.

\textbf{Equations to show:}
\begin{equation*}
\{C_j(x,y,t),\; T(x,y,t),\; \phi(x,y,t),\; n(x,y,t),\; p(x,y,t)\}
\quad \text{solved iteratively to consistency.}
\end{equation*}
\begin{align*}
D_j &= D_j(T),\\
R_j &= R_j(T,E,C_1,C_2,\dots),\\
J_n,J_p &\text{ depend on } T,\phi,\text{ and defect-assisted recombination.}
\end{align*}

\textbf{Figure idea.} A coupling loop diagram among Modules 2, 3, 4, and 5.

\textbf{Physics drawing idea.} Circular arrows connecting defect map, temperature map, potential map, and carrier map over the same device cross section.

\textbf{Speaker emphasis.} Explain that the project becomes a true degradation simulator only when these feedbacks are allowed to interact rather than being treated as one-way post-processing.

\section{Slide 8: Module 7 --- scaling, surrogates, and practical prediction}
\begin{physicsbox}
\textbf{Main question.} How does the framework move beyond one small high-fidelity simulation into useful prediction?
\end{physicsbox}

\textbf{Core message.} Module 7 is the scale-bridging layer: it uses reduced-order models, surrogate models, or statistical upscaling to extend the physically grounded simulations toward larger devices, longer times, or rarer events.

\textbf{Equations to show:}
\begin{equation*}
\mathcal{M}_{\mathrm{high\,fidelity}}
\;\longrightarrow\;
\mathcal{M}_{\mathrm{reduced/surrogate/upscaled}}
\;\longrightarrow\;
\text{practical prediction with quantified uncertainty}.
\end{equation*}

\textbf{Figure idea.} A high-fidelity local simulation feeding into a reduced or surrogate model that predicts behavior across a much larger device or many-device ensemble.

\textbf{Physics drawing idea.} Small detailed damaged-region simulation on the left, large-area chip or system-level prediction on the right, connected by an upscaling arrow.

\textbf{Speaker emphasis.} The project is scientifically valuable only if it can eventually inform prediction, not just produce beautiful local field maps. Module 7 is the bridge from mechanism to practical use.

\section{Part II: Module 4 deep-dive slides (10 slides)}

\section{Slide 9: Why ordinary Fourier heat conduction is not enough here}
\begin{physicsbox}
\textbf{Main question.} Why does Module 4 need more than the ordinary heat equation?
\end{physicsbox}

\textbf{Core message.} Radiation-related heating can be highly localized and transient. When carrier flight distances and system scales are comparable, a purely local constitutive law loses important physics.

\textbf{Equations to show:}
\begin{align*}
\qvec&=-k\nabla T,\\
C\frac{\partial T}{\partial t}&=\nabla\cdot(k\nabla T)+Q.
\end{align*}

\textbf{Figure idea.} A localized heat pulse with two conceptual predictions: an immediate smooth Fourier response versus a finite-flight-time carrier picture.

\textbf{Physics drawing idea.} Draw heat carriers leaving a small hot region; some scatter quickly while others travel ballistically before depositing energy elsewhere.

\textbf{Speaker emphasis.} The point is not that Fourier is useless. The point is that it assumes local equilibrium and local closure from the start, which can be too crude for early-time or short-length-scale transport.

\section{Slide 10: Heat carriers, mean free path, and the meaning of temperature}
\begin{physicsbox}
\textbf{Main question.} What is actually transporting the energy, and when is temperature a subtle quantity?
\end{physicsbox}

\textbf{Core message.} Heat is carried by microscopic excitations that move a finite mean free path before scattering. Outside local equilibrium, it is often better to think in terms of equivalent internal-energy temperature rather than strict thermodynamic temperature.

\textbf{Equations to show:}
\begin{align*}
\mfp &= v\tau,\\
\mathrm{Kn} &= \frac{\mfp}{L},\\
\Delta u &= C\,\Delta T.
\end{align*}

\textbf{Figure idea.} Large-device versus small-device scattering picture, plus a boundary layer with a directional nonequilibrium distribution.

\textbf{Physics drawing idea.} One inset with many randomizing collisions, one inset with nearly straight carrier flights, and a separate sketch comparing isotropic and directional carrier distributions.

\textbf{Speaker emphasis.} This slide sets up the organizing parameter for the module: the Knudsen number tells whether the thermal problem behaves diffusively, quasi-ballistically, or strongly ballistically.

\section{Slide 11: Microscopic starting point --- the Boltzmann transport equation}
\begin{physicsbox}
\textbf{Main question.} What microscopic transport equation motivates the ballistic-diffusive model?
\end{physicsbox}

\textbf{Core message.} The correct starting point is the carrier distribution function in phase space. Streaming moves carriers through space; scattering relaxes the distribution toward local equilibrium.

\textbf{Equations to show:}
\begin{equation*}
\frac{\partial f}{\partial t}+\bm{v}\cdot\nabla_{\rvec} f=-\frac{f-f_0}{\tau}.
\end{equation*}

\textbf{Figure idea.} A packet of heat carriers moving through the material while collisions gradually isotropize the distribution.

\textbf{Physics drawing idea.} Draw a local distribution function at one point deforming from a directional shape toward an isotropic equilibrium shape over a relaxation time.

\textbf{Speaker emphasis.} This is the deepest physics slide in the Module 4 sequence. Make clear that the reduced model is not ad hoc; it is a computationally cheaper descendant of the BTE.

\section{Slide 12: Diffusive limit and why Cattaneo is still incomplete}
\begin{physicsbox}
\textbf{Main question.} If Fourier fails, can Cattaneo solve the problem by itself?
\end{physicsbox}

\textbf{Core message.} Fourier conduction is the strong-scattering limit of the BTE. Cattaneo improves on Fourier by introducing finite flux relaxation, but it still treats all heat transport as one local continuum field and misses boundary-originating ballistic memory.

\textbf{Equations to show:}
\begin{align*}
\qvec&=-k\nabla T,\\
\tau\frac{\partial \qvec}{\partial t}+\qvec&=-k\nabla T.
\end{align*}

\textbf{Figure idea.} A comparison between Fourier's immediate-response picture, a Cattaneo thermal-wave picture, and a ballistic-arrival picture from a boundary.

\textbf{Physics drawing idea.} Draw one boundary-emitted carrier ray that retains memory of where it came from; contrast that with a smooth continuum wavefront.

\textbf{Speaker emphasis.} Say explicitly that finite propagation speed is necessary but not sufficient. Cattaneo still lacks a separate channel for boundary-emitted carriers.

\section{Slide 13: The ballistic-diffusive split}
\begin{physicsbox}
\textbf{Main question.} How do we keep the essential nonlocal BTE physics without solving the full phase-space problem?
\end{physicsbox}

\textbf{Core message.} Split the thermal carrier population into two parts: a ballistic piece that remembers boundaries and a medium/scattered piece that can be represented by lower-order angular moments.

\textbf{Equations to show:}
\begin{equation*}
I_\omega = I_{b\omega}+I_{m\omega}.
\end{equation*}

\textbf{Figure idea.} A material domain with straight boundary-originating rays for \(I_{b\omega}\) and a cloud-like scattered background for \(I_{m\omega}\).

\textbf{Physics drawing idea.} Color the ballistic carriers differently from the scattered carriers to show that the model is physically splitting the transport channels, not just algebraically rearranging the same heat equation.

\textbf{Speaker emphasis.} This is the conceptual heart of the module. Once the audience accepts this split, the later reduced PDE makes physical sense.

\section{Slide 14: Ballistic component by characteristics and retarded boundary time}
\begin{physicsbox}
\textbf{Main question.} How does the ballistic part carry nonlocal boundary memory into the interior?
\end{physicsbox}

\textbf{Core message.} A ballistic carrier observed at an interior point must have left the boundary earlier and survived scattering along its path. The interior ballistic state therefore depends on retarded boundary information.

\textbf{Equations to show:}
\begin{equation*}
I_{b\omega}(t,\rvec,\hat{\Omega})
=I_{w\omega}\left(t-\frac{s-s_0}{v},\rvec-(s-s_0)\hat{\Omega}\right)
\exp\left[-\int_{s_0}^{s}\frac{ds'}{v\tau_\omega}\right].
\end{equation*}

\textbf{Figure idea.} A boundary-to-interior ray with travel distance, arrival time, and exponential attenuation annotated.

\textbf{Physics drawing idea.} Draw several rays at different angles reaching one interior point; some arrive in time, some do not, which visually explains finite arrival-time causality.

\textbf{Speaker emphasis.} Emphasize that this is where the model keeps the \emph{memory} that Fourier and Cattaneo wash out. The interior field depends on what the boundary did earlier, not just on local gradients now.

\section{Slide 15: Medium/scattered component and P1 closure}
\begin{physicsbox}
\textbf{Main question.} Once the ballistic part is separated out, how is the scattered part simplified enough to solve?
\end{physicsbox}

\textbf{Core message.} The medium/scattered population is more isotropic than the ballistic part. A first-order angular expansion keeps the dominant isotropic term and the first directional correction.

\textbf{Equations to show:}
\begin{equation*}
I_{m\omega}=J_{0\omega}+\bm{J}_{1\omega}\cdot\hat{\Omega}.
\end{equation*}

\textbf{Figure idea.} Angular distribution plot showing a nearly isotropic scattered population compared with the strongly directional ballistic population.

\textbf{Physics drawing idea.} Use a polar sketch: ballistic energy concentrated into narrow angular lobes, scattered energy distributed almost uniformly with a weak directional bias.

\textbf{Speaker emphasis.} This is the step that makes the model numerically executable. The full angular dependence is not solved directly; it is reduced to spatial fields through a controlled closure.

\section{Slide 16: Final reduced ballistic-diffusive PDE}
\begin{physicsbox}
\textbf{Main question.} What is the actual reduced equation that Module 4 wants to solve?
\end{physicsbox}

\textbf{Core message.} The medium/scattered energy satisfies a Cattaneo-like PDE, but it is driven by the divergence of the ballistic heat flux. That forcing term is the key physics missing from ordinary local models.

\textbf{Equations to show:}
\begin{equation*}
\tau\frac{\partial^2 u_m}{\partial t^2}+\frac{\partial u_m}{\partial t}
=\nabla\cdot\left(\frac{k}{C}\nabla u_m\right)-\nabla\cdot\qvec_b
+\left(\dot{q}_e+\tau\frac{\partial \dot{q}_e}{\partial t}\right).
\end{equation*}

\textbf{Figure idea.} A block diagram showing \(\qvec_b\) computed from ballistic boundary transport, then fed into a PDE solve for \(u_m\), with total energy assembled from ballistic and medium pieces.

\textbf{Physics drawing idea.} One panel for ballistic rays, one panel for the smoother medium field, and one panel showing the total temperature or energy field as the sum of both contributions.

\textbf{Speaker emphasis.} Do not rush this slide. Go term by term: relaxation term, medium diffusion term, ballistic forcing term, and volumetric source term.

\section{Slide 17: Revised Module 4 target for \texttt{RadiationEffects\_proj2}}
\begin{physicsbox}
\textbf{Main question.} How is the original ballistic-diffusive derivation adapted into the new 2D Module 4 for this project?
\end{physicsbox}

\textbf{Core message.} The project does not directly solve the full angular 3D problem. It uses a 2D reduced ballistic-diffusive thermal model in silicon, with volumetric radiation-related source terms and a reduced ballistic closure appropriate for a rectangular computational domain.

\textbf{Equations to show:}
\begin{equation*}
\tau\frac{\partial^2 T_m}{\partial t^2}+\frac{\partial T_m}{\partial t}
=\nabla\cdot\bigl(\alpha\nabla T_m\bigr)
-\frac{1}{\rho_m c_p}\nabla\cdot\qvec_b
+\frac{1}{\rho_m c_p}\left(Q+\tau\frac{\partial Q}{\partial t}\right).
\end{equation*}
\begin{align*}
D_j &= D_j(T),\\
R_j &= R_j(T,E,C_1,C_2,\dots),\\
J_n, J_p &\text{ later depend on } T,\phi,\text{ and defect-assisted recombination.}
\end{align*}

\textbf{Figure idea.} A 2D silicon rectangle with localized internal heating, boundary-memory arrows, and coupling arrows to Module 3 and Module 5.

\textbf{Physics drawing idea.} Show a radiation-generated hot spot inside the domain, outward ballistic carrier paths, scattered thermal spreading, and side annotations: ``updates defect kinetics'' and ``updates carrier transport coefficients.''

\textbf{Speaker emphasis.} This is the project-specific landing slide. Explain clearly that Module 4 is now the thermal bridge between the defect-evolution physics and the later electrical-performance physics.

\section{Slide 18: Validation strategy and the plots that should matter}
\begin{physicsbox}
\textbf{Main question.} How should Module 4 be validated before it is trusted inside the larger multiphysics chain?
\end{physicsbox}

\textbf{Core message.} Validation should start with simple canonical tests such as uniform equilibrium, localized pulse spreading, and boundary heating, while also keeping the original thin-film logic in mind as a derivational benchmark. The important plots are the ones that distinguish ballistic-diffusive physics from ordinary local conduction.

\textbf{Equations to show:}
\begin{equation*}
\frac{\partial^2 \theta_m}{\partial t^{*2}}+\frac{\partial \theta_m}{\partial t^*}
=\frac{\mathrm{Kn}^2}{3}\frac{\partial^2 \theta_m}{\partial \eta^2}
-\mathrm{Kn}\frac{\partial q_b^*}{\partial \eta}.
\end{equation*}

\textbf{Figure idea.} A validation slide with four boxes: thin-film benchmark logic, 2D uniform equilibrium, 2D localized pulse, and 2D boundary heating; beside them, expected output plots such as total temperature, medium temperature, ballistic flux, and comparison against Fourier or Cattaneo baselines.

\textbf{Physics drawing idea.} For the thin film, draw finite-arrival directions reaching an interior point. For the 2D case, show a hot region whose spread is shaped by both ballistic memory and scattering.

\textbf{Speaker emphasis.} Tell the audience exactly what would count as success: equilibrium stays uniform, localized pulse spreading is finite and physically sensible, boundary heating propagates inward with retained memory, and the model trends toward Fourier behavior in the strong-scattering limit.

\section{Suggested presentation compression}
If the final talk needs to be shorter, a good compressed sequence is:
\begin{enumerate}[leftmargin=1.6em]
    \item Slide 1: project overview,
    \item one selected module slide from Part I that best matches the audience,
    \item Slide 9: why Fourier is not enough,
    \item Slide 11: BTE starting point,
    \item Slide 13 or 16: ballistic-diffusive split and final reduced PDE,
    \item Slide 17: revised 2D Module 4 in \texttt{RadiationEffects\_proj2},
    \item Slide 18: validation and expected plots.
\end{enumerate}

\end{document}
